Análisis
La distribución del saldo por segmento de clientes muestra una composición relativamente equilibrada entre los tres grupos analizados. El segmento
Premium
concentra el mayor volumen de saldos, con aproximadamente
Q8.58 millones
, seguido por el segmento
Joven
con
Q8.34 millones
y, finalmente, el segmento
Adulto
, con alrededor de
Q8.01 millones
. La diferencia entre el segmento con mayor y menor saldo es cercana a
Q500 mil
, lo que evidencia una variación moderada entre ellos.
En términos generales, no se observa una concentración marcada de los saldos en un único segmento de clientes. Los resultados reflejan una distribución relativamente homogénea, lo que sugiere que el banco mantiene una cartera diversificada y que su volumen de recursos administrados no depende exclusivamente de un perfil específico de clientes.
Aunque el segmento
Premium
registra el mayor saldo agregado, la diferencia respecto a los demás segmentos es reducida. Este comportamiento podría indicar que las estrategias de captación y retención de clientes han logrado una distribución equilibrada de los recursos entre los distintos perfiles. Sin embargo, para comprender las causas de esta distribución sería necesario complementar el análisis con información adicional, como el número de clientes por segmento, el saldo promedio por cliente y la evolución temporal de cada grupo.